\lecture{Лекция 2}{24.09.2026}

\subsection[Классификация определяющих уравнений]{Классификация определяющих уравнений для процесса формоизменения}
\label{ssec:classification}

Тензор напряжений разложен на шаровую часть и девиатор:
\[
  \sig = \underbrace{\bar\sigma\,\Eten}_{\text{шар.}} + \underbrace{\Sdev}_{\text{дев.}} ,
  \qquad
  \bar\sigma = K\bigl(\theta - 3\alpha(T-T_0)\bigr)
  \quad \text{— при тепловой деформации} .
\]
За \textbf{формоизменение} отвечает девиатор, и классификация материалов — это
классификация связи $\Sdev \leftrightarrow \dev$.

\subsubsection*{1) Упругий материал}

\[
  \Sdev = 2\mu\,\dev .
\]
Линейный пример — чистый сдвиг:
\[
  \Sdev = \tau_{12}\,(\vc i_1 \vc i_2 + \vc i_2 \vc i_1),
  \qquad \bar\sigma = 0
  \quad\Longrightarrow\quad
  \tau_{12} = 2\mu\,\dev_{12},
\]
$\mu$ — модуль сдвига. Из II и III законов термодинамики следует $\mu \ge 0$.

\begin{warn}
Запись чистого сдвига ($\Sdev$ через диады $\vc i_1 \vc i_2 + \vc i_2 \vc i_1$ и
пометка «$\sigma' = 0$») на фото неразборчива — восстановлена по смыслу.
\end{warn}

Реологический элемент закона Гука — \textbf{пружина}.

\inputfig{figures/fig-rheo-elements}

\subsubsection*{2) Вязкий материал (жидкость)}

\[
  \Sdev = 2\eta\,\dt{\dev} .
\]
Из II закона термодинамики $\eta \ge 0$. Реологический элемент Ньютона —
\textbf{демпфер}.

\subsubsection*{3) Пластический материал}

При недостижении некоторого предела деформации нет.

\inputfig{figures/fig-plastic-diagram}

\[
  N(\Sdev) < \tau_s \quad\Longrightarrow\quad \dt{\dev} = 0
  \quad(\text{скорость деформации} = 0),
\]
\[
  N(\Sdev) = \tau_s \quad\Longrightarrow\quad \Sdev = \lambda\,\dt{\dev} .
\]
Напряжение $\Sdev$ не может расти вместе с ростом $\dt{\dev}$, поэтому
подстраивается множитель $\lambda$. Из термодинамики $\lambda \ge 0$.

Реологический элемент Сен-Венана (сухое трение) — это такой материал.

\begin{warn}
$N(\Sdev)$ — норма девиатора (интенсивность касательных напряжений),
$\tau_s$ — предел текучести; явные формулы на лекции не записаны.
\end{warn}
